\scanpage{28}
Alors la suite spectrale implique
\[
H^i_Y(X,\mu_{nX}^{\otimes d})=0\qquad\text{si }i<2d
\]
\[
H^{2d}_Y(X,\mu_{nX}^{\otimes d})\simeq\Gamma(Y,\underline H^{2d}_Y(\deleted{\unsure{\mathbf Z/n\mathbf Z}}\mu_{nX}^{\otimes d}))
\]
d'où un \deleted{\unsure{hom} canonique} élément canonique
\[
\gamma_{X/Y}^{(n)}\in H^{2d}_Y(X,\mu_{nX}^{\otimes d})
\]
[qu'on peut maintenant spécialiser dans $H^{2d}_{\deleted{\unsure{Y}}}(X,\mu_{nX}^{\otimes d})$, ou si $X$ est propre sur un $S$ dans
\[
H^{2d}_{\Phi_{Y/X}}(X,\mu_{nX}^{\otimes d}).
\]
\margin{$\Phi_{Y/X}\ \longrightarrow\ \text{famille de \unsure{cpts}}$}

\margin{N.B ($\beta_{Y/X}\in H^1_Y(X,G_{mX})$ et $\gamma^{(n)}_{Y/X}=\partial_n(\beta_{Y/X})$, en particulier son image dans $H^2_Y(X,G_{mX})$ est nulle)}

\uline{Exemple 1}. $X$ lisse sur $k$ \inserted{sép. clos.}, \deleted{$x$ pt} dim relative $d$, \deleted{\unsure{il va résulter des} $\unsure{C}$ \unsure{th. de}} \inserted{$x$ pt de $X$ \unsure{rat} sur $k$, d'où}
\[
\gamma_{x/X}\in H^{2d}_!(X,\mu_{nX}^{\otimes d}).
\]
Il va résulter de le théor. qu'on va \unsure{prouver} que \deleted{\unsure{pour} $\unsure{A}$} si $X$ est géom. connexe, cet élément est \uline{une base} \uline{de} $H^{2d}_!(X,\mu_{nX}^{\otimes d})$ \uline{et ne dépend pas du choix de $x$}; \deleted{\unsure{ici}} C'est qu'on avait vérifié directement dans le cas où $d=1$, en \deleted{\unsure{le compa}} \inserted{nous \unsure{ramenant}} au cas où $C$ étant complet et où on dispose déjà
\mspage{27}
